% =====================================================================
%  Chapter 20 — The Proof in Natural Language, 4
% =====================================================================

\chapter{The Determinant and the Interlacing Weights}
\label{ch:determinant}

The arithmetic factor of the previous chapter produces a count. What the
argument needs is a matrix.

\section{From a sum to a determinant}Recall the goal: exhibit a matrix $M(g)$ whose determinant is the Weil
expression and which is a Gram matrix. The natural route is to note that a
determinant is a signed sum over permutations,
\begin{equation}
\det M = \sum_{\sigma \in \mathfrak{S}_n} \operatorname{sgn}(\sigma)
\prod_{i} M_{i,\sigma(i)},
\end{equation}
and that a Gram matrix's determinant is a signed sum of products of inner
products --- so that if the entries are inner products of vectors $v_i$, the
determinant is the squared volume of the parallelepiped they span, and is
non-negative automatically.

\begin{quote}
\textbf{Proposition I (detector).} Suppose the Weil expression admits the
expansion
\[
W(T) \;=\; \sum_{\rho}\sum_{\sigma} \operatorname{sgn}(\sigma)
\prod_{i} \hat g(\rho_i)\,\hat g(1-\rho_{\sigma(i)}),
\]
with the indexing arranged so that the terms with equal sign carry
\emph{distinct} arguments. Then $W(T) \ge 0$.
\end{quote}

The content of the proposition is the last clause. A signed sum is
non-negative if the positive terms dominate the negative ones; it is
non-negative for a structural reason if the terms with opposite signs cannot
pair up. The remark's phrase is exactly this: ``opposite signs meaning no
cancellation of inverse roots''.

\section{The interlacing weights}Where do the signs come from? From the interlacing of the arguments of the
zeros, and of the periods.

\begin{description}[style=nextline,leftmargin=0pt,labelwidth=0pt,labelsep=0pt,font=\sffamily\bfseries\small,
  itemsep=6pt]
\item[Zeros.] The non-trivial zeros of $\zeta$ have arguments that interlace:
as the ordinates increase, the arguments of the zeros alternate between the
two sides of the critical line.
\item[Periods.] The Gross--Prasad periods attached to the Hermitian forms
have arguments of their own, and the parameter data of the argument records
where they sit relative to the arguments of the zeros.
\end{description}

The remark's item 7:

\begin{quote}
\small
``Determinants for $W(\cdot)$'s containing the desired periods, indeed GGP
periods, with values 0, 1 meaning interlacing weights, or up to $m$
non-interlacing ones.''
\end{quote}

So the entries of the matrix are the GGP periods; the values $0$ and $1$ are
the \emph{interlacing weights} --- a weight of $1$ when the period's argument
interlaces the zero's and $0$ when it does not --- and the statement allows up
to $m$ non-interlacing entries, where $m$ is the length of the chain of
lifts.

\begin{quote}
\textbf{Proposition J.} The interlacing weights, together with the separation
condition of Chapter~\ref{ch:count}, give the sign structure required by
Proposition~I, and hence $W(T) \ge 0$.
\end{quote}

\section{Why the interlacing must be manufactured}Here is the difficulty that the function-field proof does not have, and that
this method has to solve.

In the classical Weil bound, the interlacing is \emph{given}: it is a theorem
about the arguments of the Frobenius eigenvalues on the Jacobian of a curve,
and it is proved by Deligne's theory of the Weil conjectures or, in the
elementary case, by the geometry of the Jacobian. The corresponding
eigenvalues are those of an actual endomorphism, and their arguments
interlace because the endomorphism preserves the polarisation.

Here there is no curve, no Jacobian, and no endomorphism. The data are
automorphic periods, and the interlacing has to be imposed or produced. The
formalisation carries the interlacing as part of the parameter data
(\code{P.Interlace} in the \code{Hyp3} condition of
\code{FunctionModelRallisFromInputs.lean}), together with the parity and
centrality conditions.

\begin{quote}
\textbf{Proposition K (the gap).} The parameter data must satisfy the
interlacing condition, the parity condition, the centrality condition, and
the cubic separation condition of the previous chapter; and the existence of
such data, with the arithmetic properties the argument requires, is the
non-trivial content.
\end{quote}

\section{The formal statements}The development's machinery for this is substantial. The files that carry it:

\begin{center}
\begin{tabular}{@{}>{\RaggedRight\arraybackslash\hyphenpenalty=10000\exhyphenpenalty=10000}p{0.44\textwidth}>{\RaggedRight\arraybackslash\hyphenpenalty=10000\exhyphenpenalty=10000}p{0.44\textwidth}@{}}
\toprule
\textbf{File} & \textbf{What it establishes} \\
\midrule
\code{.../DecoratedLatticeProofs/WeilCoefficientPureTensorEqSmulInnerStatement.lean} &
the Weil coefficient is an inner product of a conjugate automorphic form with a
cuspidal one \\
\code{.../WeilCoefficientExternalProductEqMulOrthogonalSumStatement.lean} &
external products of Weil coefficients multiply, with orthogonal sums \\
\code{.../WeilCoefficientSplitsLocalFactorTwistPlaceStatement.lean} &
the coefficient splits into local factors at each place \\
\code{.../ThetaInnerProductFormulaDefiniteOfRallisConverge.lean} &
the theta inner product formula, and the convergence of the Rallis
construction \\
\code{.../FrameEstimates/FrameBasis.lean} &
the frame basis whose Gram matrix is the object being constructed \\
\code{.../FrameEstimates/ExistsColsDetNeZero.lean} &
the non-vanishing of the relevant determinant \\
\midrule
\end{tabular}
\end{center}

\section{The three parts of the positivity}It is worth separating what the determinant argument needs, because the three
are usually conflated.

\begin{enumerate}
\item \textbf{The Gram structure.} The matrix must be a Gram matrix: $M_{ij}
= \langle v_i, v_j\rangle$. The Weil coefficient statements above are
precisely the assertions that the entries are inner products, and that inner
products compose.
\item \textbf{The sign structure.} The interlacing must give the alternating
signs. This is what makes the determinant non-negative, and it is where the
separation condition is needed.
\item \textbf{The non-vanishing.} The determinant must not be identically
zero. \code{ExistsColsDetNeZero} is the statement that some choice of columns
gives a non-zero determinant --- without which the Gram matrix is degenerate
and the positivity is vacuous.
\end{enumerate}

\section{The honest position}Parts 1 and 3 are formalised, and the formalisation is substantial. Part 2 ---
the sign structure --- is where the interlacing condition enters as
\emph{input} rather than as a derived consequence, and whether the parameter
data satisfying it can be produced is the question.

That is not a gap in the formalisation. It is a gap in the mathematics, of the
same kind as the two in Chapter~\ref{ch:budget} and
Chapter~\ref{ch:vanishing}, and it is located here rather than at the end
because that is where it occurs.

\begin{keytakeaways}
\item The positivity of the Weil expression is a Gram-matrix statement, and
the argument constructs the Gram matrix out of GGP periods with interlacing
weights.
\item The interlacing of the arguments of the zeros --- which in the
function-field proof is a theorem about Frobenius on a Jacobian --- must here
be produced or imposed, and that is the new content.
\item Three things are needed: the Gram structure, the sign structure, and the
non-vanishing of the determinant. The first and third are formalised; the
second enters as a condition on the parameter data.
\end{keytakeaways}
